% GENERATED FILE. Edit canonical lessons, then run npm run generate:books.
% canonical-source-hash: fnv1a64:b910a535f0071113
% canonical-lessons: ZH-C308-shizong, ZH-C308-zicong, ZH-W308-zi-self, ZH-C308-laizi, ZH-C308-fanzui, ZH-C308-paichusuo

\chapter{Ever Since}
\label{ch:zh-ever-since}
\hlchaptermodality{308}

\begin{chapteropening}
\textbf{By the end of this chapter:} \emph{I can say since when, where someone comes from, and that a person went missing.}

Complete the last lesson of chapter 308: I can say since when, where someone comes from, and that a person went missing.
\end{chapteropening}

\section[shīzōng]{shīzōng — to go missing}

\label{lesson:ZH-C308-shizong}



\begin{quote}
Before the new one: say the Mandarin for a honeymoon, then the Mandarin for a white wedding dress.
\end{quote}

\subsection*{You'll want to know: shīzōng}
\emph{shīzōng} — \textquotedblleft{}to go missing\textquotedblright{}.

\begin{quote}
\emph{nàge háizi shīzōng le} — That child has gone missing.
\end{quote}

Its characters come later; for now it is a word for the ear.

\begin{cousinweb}
\emph{shīzōng} is built from \emph{shī} 'to lose' and \emph{zōng} 'a footprint', both inherited from Middle Chinese, \emph{syit} and \emph{tsjowng} (Guangyun, in Baxter's transcription), which glosses the first 'a mistake' and 'to let go' and the second 'tracks': the trail is lost.
\end{cousinweb}

\begin{grammarlens}[title={What you've built}]
One more word for telling what happened.
\end{grammarlens}

\subsection*{Your turn}
\begin{itemize}
\raggedright
  \item \emph{Say it:} \emph{shīzōng}
  \item \emph{Say it:} \emph{shīzōng}, once more
  \item \emph{Say it:} \emph{nàge háizi shīzōng le} — That child has gone missing.
  \item \emph{Recall:} say the Mandarin for a honeymoon, then the Mandarin for a white wedding dress, then say \emph{shīzōng} again
\end{itemize}

\begin{tcolorbox}[breakable,colback=teal!4,colframe=teal!35!black,title={Before you move on}]
What does shīzōng mean? (\textquotedblleft{}To go missing\textquotedblright{}.) Say it once more.
\end{tcolorbox}
\section[\texorpdfstring{\zh{自从} (zìcóng)}{自从 (zìcóng)}]{\zh{自从} (zìcóng) — since; ever since}

\label{lesson:ZH-C308-zicong}



\begin{quote}
Before the new one: say the Mandarin for a white wedding dress, then the Mandarin for to go missing.
\end{quote}

\subsection*{You'll want to know: \zh{自从}}
\textbf{\zh{自从}} — \emph{zìcóng} — \textquotedblleft{}since; ever since\textquotedblright{}.

\begin{quote}
\textbf{\zh{自从那天以后我天天写日记}} \emph{zìcóng nà tiān yǐhòu, wǒ tiāntiān xiě rìjì} — Since that day, I've kept a diary every day.
\end{quote}

The new character, in pieces: \zh{自} — six strokes: a short sweep down to the left on top, then a tall frame: a vertical, a bar that turns down the right side, two bars inside, and the bar that closes the foot, like \zh{日} grown taller. The next lesson puts it in your hand.

\begin{cousinweb}
\textbf{\zh{自从}} (\emph{zìcóng}) is built from \zh{自} 'from' and \zh{从} 'from', both inherited from Middle Chinese, \emph{dzijH} and \emph{dzjowng} (Guangyun, in Baxter's transcription), which glosses the first 'from' and 'by way of' and the second 'to go along with'.
\end{cousinweb}

\begin{grammarlens}[title={What you've built}]
One more word for telling what happened.
\end{grammarlens}

\subsection*{Your turn}
\begin{itemize}
\raggedright
  \item \emph{Say it:} \emph{zìcóng}
  \item \emph{Say it:} \emph{zìcóng}, once more
  \item \emph{Say it:} \emph{zìcóng nà tiān yǐhòu, wǒ tiāntiān xiě rìjì} — Since that day, I've kept a diary every day.
  \item \emph{Recall:} say the Mandarin for a white wedding dress, then the Mandarin for to go missing, then say \emph{zìcóng} again
\end{itemize}

\begin{tcolorbox}[breakable,colback=teal!4,colframe=teal!35!black,title={Before you move on}]
What does \zh{自从} mean? (\textquotedblleft{}Since; ever since\textquotedblright{}.) Say it once more.
\end{tcolorbox}
\section[\texorpdfstring{\zh{自} (zì)}{自 (zì)}]{\zh{自} — a short sweep over a tall frame}

\label{lesson:ZH-W308-zi-self}



\begin{quote}
The word first, then the character inside it.

\begin{itemize}
\raggedright
  \item \emph{Recall:} say \emph{zìcóng} — \textbf{R1}, one lesson back, before it decays
\end{itemize}

\begin{quote}
\zh{日} \zh{白}
\end{quote}

\emph{Read it:} each one aloud, with nothing said first — \emph{rì}, \emph{bái}
\end{quote}

\subsection*{Script — a short sweep over a tall frame}
\hlblockfigure{\detokenize{figures/ZH-W308-zi-self-filmstrip.pdf}}{How \zh{自} is written, stroke by stroke}

\begin{quote}
\zh{自}
\end{quote}

\textbf{Two pieces: a short sweep on top, a tall frame under it.} Count the pieces, not the strokes.

\textbf{The sweep, one stroke}: short, down to the left, at the top middle.

\textbf{The frame, five strokes}, built like \zh{日}: a vertical down the left; a bar across the top that turns down the right side; two bars inside; then the bar that closes the foot.

1 + 5 = six.

\subsection*{Your turn}
\begin{itemize}
\raggedright
  \item \emph{Write it:} \zh{自}, all six, with the model in front of you
  \item \emph{Cover:} write it again from memory, and count the pieces, not the strokes
\end{itemize}

\begin{tcolorbox}[breakable,colback=teal!4,colframe=teal!35!black,title={Before you move on}]
How many strokes in \zh{自}? (\textbf{Six}.) What comes first? (\textbf{The short sweep on top}.) How many bars inside the frame? (\textbf{Two}.)

Source: \href{https://raw.githubusercontent.com/chanind/hanzi-writer-data/68d10a4b21150cae5e1ebbd223eed289cf32d90c/data/\%E8\%87\%AA.json}{Hanzi Writer Data \zh{自}, pinned stroke paths and medians}.
\end{tcolorbox}
\section[\texorpdfstring{\zh{来自} (láizì)}{来自 (láizì)}]{\zh{来自} (láizì) — to come from}

\label{lesson:ZH-C308-laizi}



\begin{quote}
Before the new one: say the Mandarin for to go missing, then the Mandarin for ever since.
\end{quote}

\subsection*{You'll want to know: \zh{来自}}
\textbf{\zh{来自}} — \emph{láizì} — \textquotedblleft{}to come from\textquotedblright{}.

\begin{quote}
\textbf{\zh{来自中国的新生经常来我家}} \emph{láizì Zhōngguó de xīnshēng jīngcháng lái wǒ jiā} — The new students from China often come to our house.
\end{quote}

\begin{cousinweb}
\textbf{\zh{来自}} (\emph{láizì}) is built from \zh{来} 'to come' and \zh{自} 'from', both inherited from Middle Chinese, \emph{loj} and \emph{dzijH} (Guangyun, in Baxter's transcription), which glosses the first 'to arrive'.
\end{cousinweb}

\begin{grammarlens}[title={What you've built}]
One more word for telling what happened.
\end{grammarlens}

\subsection*{Your turn}
\begin{itemize}
\raggedright
  \item \emph{Say it:} \emph{láizì}
  \item \emph{Say it:} \emph{láizì}, once more
  \item \emph{Say it:} \emph{láizì Zhōngguó de xīnshēng jīngcháng lái wǒ jiā} — The new students from China often come to our house.
  \item \emph{Recall:} say the Mandarin for to go missing, then the Mandarin for ever since, then say \emph{láizì} again
\end{itemize}

\begin{tcolorbox}[breakable,colback=teal!4,colframe=teal!35!black,title={Before you move on}]
What does \zh{来自} mean? (\textquotedblleft{}To come from\textquotedblright{}.) Say it once more.
\end{tcolorbox}
\section[fànzuì]{fànzuì — to commit a crime}

\label{lesson:ZH-C308-fanzui}



\begin{quote}
Before the new one: say the Mandarin for ever since, then the Mandarin for to come from a place.
\end{quote}

\subsection*{You'll want to know: fànzuì}
\emph{fànzuì} — \textquotedblleft{}to commit a crime\textquotedblright{}.

\begin{quote}
\emph{tā shuō tā méiyǒu fànzuì} — He says he hasn't committed any crime.
\end{quote}

Its characters come later; for now it is a word for the ear.

\begin{cousinweb}
\emph{fànzuì} is built from \emph{fàn} 'to violate' and \emph{zuì} 'a crime', both inherited from Middle Chinese, \emph{bjomX} and \emph{dzwojX} (Guangyun, in Baxter's transcription), which glosses the first 'to trespass' and 'to violate' and records that the second was once written with another character.
\end{cousinweb}

\begin{grammarlens}[title={What you've built}]
One more word for telling what happened.
\end{grammarlens}

\subsection*{Your turn}
\begin{itemize}
\raggedright
  \item \emph{Say it:} \emph{fànzuì}
  \item \emph{Say it:} \emph{fànzuì}, once more
  \item \emph{Say it:} \emph{tā shuō tā méiyǒu fànzuì} — He says he hasn't committed any crime.
  \item \emph{Recall:} say the Mandarin for ever since, then the Mandarin for to come from a place, then say \emph{fànzuì} again
\end{itemize}

\begin{tcolorbox}[breakable,colback=teal!4,colframe=teal!35!black,title={Before you move on}]
What does fànzuì mean? (\textquotedblleft{}To commit a crime\textquotedblright{}.) Say it once more.
\end{tcolorbox}
\section[pàichūsuǒ]{pàichūsuǒ — local police station}

\label{lesson:ZH-C308-paichusuo}



\begin{quote}
Before the new one: say the Mandarin for to come from a place, then the Mandarin for to commit a crime.
\end{quote}

\subsection*{You'll want to know: pàichūsuǒ}
\emph{pàichūsuǒ} — \textquotedblleft{}local police station\textquotedblright{}.

\begin{quote}
\emph{wǒmen qù pàichūsuǒ bàojǐng} — We went to the police station to report it.
\end{quote}

Its characters come later; for now it is a word for the ear.

\begin{cousinweb}
\emph{pàichūsuǒ} is built from \emph{pài}, \zh{出} 'out' and \zh{所} 'a place', all inherited from Middle Chinese, \emph{pheaH}, \emph{tsyhwit} and \emph{srjoX} (Guangyun, in Baxter's transcription), which glosses the first 'a branch of a river'; its sense 'to send' is not given there. The place that sends officers out.
\end{cousinweb}

\begin{grammarlens}[title={What you've built}]
One more word for telling what happened.
\end{grammarlens}

\subsection*{Your turn}
\begin{itemize}
\raggedright
  \item \emph{Say it:} \emph{pàichūsuǒ}
  \item \emph{Say it:} \emph{pàichūsuǒ}, once more
  \item \emph{Say it:} \emph{wǒmen qù pàichūsuǒ bàojǐng} — We went to the police station to report it.
  \item \emph{Recall:} say the Mandarin for to come from a place, then the Mandarin for to commit a crime, then say \emph{pàichūsuǒ} again
\end{itemize}

\begin{tcolorbox}[breakable,colback=teal!4,colframe=teal!35!black,title={Before you move on}]
What does pàichūsuǒ mean? (\textquotedblleft{}Local police station\textquotedblright{}.) Say it once more.
\end{tcolorbox}
